\chapter{Literature Review}\label{chap:literature_review}

\section{Post-Harvest Loss and the Role of Mould}

Strawberries are among the most commercially valuable soft fruits grown in Europe, yet they are also among the most fragile. Global production exceeded nine million tonnes in 2022 \cite{FAOSTAT2024}, and across the European Union, fresh fruit consistently ranks as one of the most wasted food categories at every stage of the supply chain \cite{EC2019}. Post-harvest decay is one of the primary causes of that loss, and among the pathogens responsible, \textit{Botrytis cinerea} is the most damaging. The grey mould it causes is present in virtually every strawberry-growing region in the world and can spread rapidly through a batch under the right conditions \cite{Petrasch2019}.

What makes \textit{B. cinerea} particularly difficult to manage in a supply chain context is its relationship with the very conditions used to preserve freshness. Strawberries are typically stored and transported at low temperatures and high relative humidity to slow ripening and moisture loss. Those same conditions, particularly when temperatures drift upward or humidity is poorly controlled, also favour mould establishment. Temperature and atmospheric composition during storage have a direct influence on how quickly grey mould develops and spreads \cite{Nunes2012}. Once visible spoilage appears on a batch, the commercial value of the entire shipment is at risk --- not just the directly affected fruit.

This creates a detection problem that post-harvest researchers have long recognised. By the time grey mould is visible to an inspector, it has already been present and metabolically active for some time. The window for effective intervention --- adjusting temperature, increasing ventilation, segregating cargo --- is narrow, and it closes before manual inspection can provide the necessary information.

\section{Environmental Monitoring in Cold Chain Logistics}

The cold chain for fresh produce is not a single controlled environment. From packhouse to retail, strawberries pass through refrigerated trucks, distribution centres, loading bays, and retail cold rooms, each with its own microclimate and each introducing the possibility of a temperature or humidity excursion. Managing conditions consistently across these transitions is a well-documented challenge: even brief departures from the recommended storage range can accelerate mould development and reduce shelf life significantly \cite{Mercier2017}.

The problem is compounded by spatial variability within a single environment. Studies deploying wireless sensor networks in commercial growing and storage facilities have recorded temperature differences of up to \SI{3.3}{\celsius} and humidity differences of up to 9\% between locations within the same enclosed space \cite{Ferentinos2017}. A single sensor positioned at one point in a storage room or truck compartment will not represent the conditions experienced by all of the produce inside it.

IoT-based sensor networks have emerged as a practical response to this variability, offering continuous spatially distributed measurements that periodic manual inspection cannot \cite{Lamberty2025}. However, scaling these systems across distributed fleets and remote facilities introduces constraints that simple monitoring architectures struggle to meet. Network connectivity cannot be assumed for vehicles in transit or rural storage sites. Cloud-dependent systems that rely on remote servers for data processing or model inference fail whenever the link is lost. Fixed-threshold alarms handle straightforward exceedances but cannot model the multi-variable nature of mould risk, where the combined behaviour of temperature, humidity, and airborne chemistry determines the actual spoilage hazard \cite{Wang2024}. The gap between what current monitoring approaches provide and what effective early mould detection requires is the starting point for the work in this thesis.

\section{Volatile Organic Compounds as Mould Indicators}

One of the most well-studied approaches to non-destructive mould detection is the measurement of volatile organic compounds (VOCs). Fungi produce a characteristic range of VOCs as metabolic byproducts throughout their growth cycle, and this production begins in the early stages of infection, before any visible decay is established \cite{Yang2025}. For \textit{Botrytis cinerea} specifically, identified biomarkers in infected strawberry include 1-octen-3-ol, 2-methyl-1-butanol, 2-methyl-1-propanol, and short-chain ketones such as 1-octen-3-one \cite{Vandendriessche2012}. Because they are released during early fungal metabolism, changes in the VOC profile of a storage environment can serve as a leading indicator of mould onset rather than a confirmation of spoilage already underway \cite{Ma2023}.

This principle underpins a substantial body of work in the electronic nose (e-nose) field. Baietto and Wilson \cite{BaiettoWilson2015} provide a broad review of e-nose applications across the fruit production chain, covering fruit identification, ripeness assessment, and quality grading. Their work demonstrates that gas sensor arrays can reliably detect and discriminate the volatile profiles associated with different fruit species and decay states, establishing the scientific basis for the sensor-based approach taken in this thesis. The limitation of most e-nose research, as they acknowledge, is that the systems involved are laboratory instruments --- expensive, bulky, and not designed for continuous unattended deployment. Translating that detection capability into a low-cost autonomous node is precisely the engineering challenge addressed here.

A practical consideration that is often underappreciated in VOC sensor deployment is the physical behaviour of the gases being measured. Mould fungi broadly produce MVOCs including ethanol, 1-octen-3-ol, and short-chain alcohols and ketones \cite{Zhao2025}, and \textit{Botrytis cinerea} is established as a producer of these same volatile classes during its metabolic activity \cite{Vandendriessche2012}. Many of these compounds have molecular weights substantially higher than that of air (approximately 29\,g/mol): ethanol is 46\,g/mol, and 1-octen-3-ol reaches 128\,g/mol. In low-airflow enclosed environments such as refrigerated trucks or cold storage units, denser gases do not disperse uniformly but stratify under gravity, accumulating near the base of the space rather than at the height of the stored produce \cite{Zhang2026}. This has a direct implication for sensor node positioning: a sensor mounted level with the fruit will sit above the primary zone of gas accumulation and will underestimate the actual VOC concentration building up below it. Accounting for this gravitational stratification is an important and often overlooked design consideration in mould detection sensor systems \cite{Suawa2025}.

\section{TinyML on Microcontrollers}

TinyML describes the class of machine learning deployments that run on microcontrollers and similarly constrained embedded devices --- hardware operating with kilobytes of memory and milliwatt-scale power budgets \cite{Dutta2021}. The field has developed rapidly as a direct response to the connectivity and energy costs of cloud-dependent inference, and as embedded frameworks have matured to support increasingly capable models within these constraints.

Three frameworks are relevant to the work in this thesis, each representing a different position on the spectrum of on-device capability. TensorFlow Lite for Microcontrollers (TF Lite Micro) \cite{David2021} converts a pre-trained model into a compact quantised representation and runs inference only. No learning occurs on the device after deployment. It is the most mature option and the most conservative in its resource demands. TinyOL \cite{Disabato2021} extends this with partial on-device adaptation: the bulk of the network is frozen, but the final layer can be updated from new observations after deployment, allowing the output mapping to shift without retraining the feature extraction layers. AIfES \cite{Wulfert2024} goes furthest, supporting complete forward and backward passes on the microcontroller itself. A network can be trained from scratch in the field, with no pre-trained weights required.

The reason this spectrum matters for cold chain monitoring comes back to the variability described in earlier sections. A model trained on data from one environment --- or on aggregated data from many environments --- is not guaranteed to perform reliably at a new deployment site with its own temperature baseline, humidity profile, and VOC background. This is concept drift: the statistical relationship between sensor inputs and mould risk labels shifts as the environment changes, and a static model will degrade silently rather than signalling that it is no longer calibrated for current conditions \cite{Bayram2022, Ilic2023}. A node that can retrain locally on data from its own environment is in a fundamentally stronger position. The open question --- and the one this thesis addresses --- is what that capability costs in energy terms, and whether the cost is operationally justified for a battery-powered autonomous deployment.

On-device training is not a new idea, but systematic energy characterisation of the available frameworks on commodity microcontroller hardware remains limited. Federated approaches that distribute training across networks of devices have attracted significant research attention \cite{Ficco2024}, but these still assume periodic connectivity for model aggregation. The fully autonomous case --- a single disconnected node training and adapting without any external coordination --- has received less attention in terms of concrete energy measurement, which is the gap this thesis fills.

\section{Energy Profiling for Embedded Systems}

Energy consumption is the defining constraint on autonomous embedded deployments. A sensor node that classifies mould risk accurately but exhausts a battery within days provides no practical advantage over a manual inspection regime. Understanding and quantifying energy use across all operational phases is therefore not a secondary concern --- it is central to evaluating whether a given framework is viable for the intended deployment.

For microcontroller-based systems, energy profiling is typically performed by measuring current at the supply rail during distinct operational phases and combining these measurements with timing data to derive energy per operation. The Nordic Semiconductor Power Profiler Kit II (PPK2) \cite{PPK2datasheet} has become a standard instrument for this kind of work, functioning simultaneously as a regulated power supply and a high-resolution current measurement device. It is capable of capturing transient current events at the microsecond timescale, making it suitable for resolving both the brief high-current bursts associated with computation-intensive operations and the low quiescent current of deep sleep modes.

For a duty-cycled node --- one that wakes on a schedule to sample sensors, runs any necessary inference or training, and returns to sleep --- battery lifetime can be estimated analytically from measured current consumption figures without requiring live discharge testing, an approach with established precedent in wireless sensor network research \cite{Rodrigues2017}. This approach allows different operational configurations to be compared in a controlled way: varying training frequency, sleep duty cycle, or sensor activation patterns can all be modelled from the same set of measurements. It also separates the energy cost of the ML workload from the baseline cost of sensing, which is a necessary distinction when the goal is to understand what on-device training specifically adds to the power budget.

\section{Summary}

The literature traces a clear line of reasoning from the agricultural problem to the technical approach. Strawberry post-harvest loss driven by \textit{Botrytis cinerea} is a commercially significant problem whose severity is directly tied to environmental conditions --- conditions that are, in principle, continuously monitorable. VOCs produced during early fungal metabolism offer a non-destructive early indicator of mould onset, and the physical behaviour of those compounds informs where sensors should be placed relative to the stored produce. Existing monitoring approaches --- manual inspection, fixed-threshold alarms, and cloud-connected machine learning --- each fall short under the constraints of a distributed, disconnected cold chain deployment, and the microclimate variability between sites means that a centrally trained model cannot be expected to transfer reliably. TinyML frameworks address these limitations by bringing adaptive inference and, in some cases, full on-device training to embedded hardware. But the frameworks differ substantially in what they require from the device, and the energy cost of that difference has not been thoroughly characterised.

\textbf{Gap in the literature.} Three specific gaps motivate this work. First, while TinyML frameworks such as TF Lite Micro, TinyOL, and AIfES have each been demonstrated individually on embedded hardware, no study has placed all three in direct empirical comparison on the same device, dataset, and measurement infrastructure. Existing comparisons are either theoretical, limited to inference-only configurations, or evaluated on platforms that differ in architecture from the commodity ESP32 hardware prevalent in low-cost IoT deployments. Second, the energy cost of full on-device neural network training --- as opposed to inference or output-layer adaptation --- has not been measured at the per-operation level on a microcontroller in an agricultural monitoring context. Federated and distributed learning approaches have received attention \cite{Ficco2024}, but these still assume periodic connectivity; the fully autonomous case, where a single disconnected node trains and adapts without any external coordination, remains uncharacterised in terms of concrete energy figures. Third, no prior work has evaluated the battery lifetime implications of framework selection for a multi-sensor mould detection node under a realistic duty cycle, making it impossible for practitioners to judge whether full on-device training is operationally viable for battery-powered deployments. This thesis addresses all three gaps through direct PPK2 measurement of each framework on an ESP32 sensing node, producing the per-operation energy figures, accuracy results, and battery lifetime estimates needed to make that judgement concrete.

The following chapters present the system design, experimental methodology, and results that quantify that trade-off and assess its operational implications.
